% ============================================================
% INTRODUCTION GÉNÉRALE (1 page pleine, paragraphes)
% ============================================================
\chapter*{Introduction Générale}
\addcontentsline{toc}{chapter}{Introduction Générale}
\markboth{INTRODUCTION GÉNÉRALE}{}
\thispagestyle{plain}

\begingroup
\setstretch{1.22}
\setlength{\parskip}{0.55em}

Aujourd'hui, le recrutement est un enjeu majeur pour les entreprises. Il ne
suffit plus de publier une offre. Il faut aussi attirer les bons profils, traiter
rapidement les candidatures et suivre clairement chaque étape, de l'annonce
jusqu'à l'embauche. La qualité du recrutement influence directement la performance
des équipes, la stabilité des effectifs et l'image de l'entreprise auprès des
candidats.

Dans un réseau multi-agences comme \textbf{DAAM}, ce besoin est encore plus fort.
Plusieurs responsables RH travaillent sur des zones différentes. Ils doivent rester
organisés, éviter les doublons et ne pas perdre la vision d'ensemble. Sans un outil
partagé, la coordination devient difficile. Le suivi manque alors d'homogénéité
d'une agence à l'autre, et la direction dispose rarement d'une information fiable
en temps utile.

Sans outil unique, les informations se dispersent facilement entre fichiers, e-mails
et suivi manuel. Les décisions deviennent plus lentes, le parcours des candidats
manque de clarté, et l'expérience des postulants peut se dégrader. Une solution
centralisée devient alors nécessaire pour gagner en efficacité, en cohérence et
en transparence, tout en gardant un suivi simple et lisible pour les équipes RH.

C'est dans ce contexte que s'inscrit ce \textbf{Projet de Fin d'Année (PFA)},
réalisé pendant environ \textbf{deux mois} au sein de \textbf{DAAM}. L'objectif
est de concevoir et développer une \textbf{plateforme de gestion du recrutement}.
Cette plateforme doit digitaliser le parcours complet, de la création des offres
jusqu'à la confirmation d'embauche. Elle doit aussi centraliser le travail des
rôles Administrateur, RH et Candidat, avec une vue globale pour l'administration.
Enfin, elle doit aider à une meilleure sélection grâce au suivi des statuts,
aux tests et aux tableaux de bord.

Le projet a été mené selon une approche agile SCRUM, avec une avancée progressive
des fonctionnalités et une validation régulière des besoins. Cette méthode a permis
d'organiser le travail de façon claire et d'adapter la solution aux attentes des
utilisateurs. Au fil des sprints, la plateforme a intégré notamment un module de
\textbf{réclamations} et une chaîne \textbf{DevOps} (Docker, Jenkins, SonarQube).
Des perspectives demeurent, comme l'enrichissement des notifications in-app.

Ce rapport présente d'abord un état des lieux du recrutement dans le chapitre~1.
Le chapitre~2 situe ensuite le projet dans son contexte chez DAAM. Les chapitres
suivants détaillent l'analyse des besoins, la conception, la réalisation par
sprints, puis l'industrialisation DevOps. L'ensemble vise à montrer comment un
outil unique et bien organisé peut améliorer le quotidien des RH et le suivi
global du recrutement.

\endgroup
\clearpage
